%%%% Phase-2 technical report for team Antarctic penguins.
%%%% TPC @ IJCAI-ECAI 2026 Phase 2 technical report.
%%%% Uses the unmodified official IJCAI-ECAI 2026 author kit.

\typeout{TPC 2026 Phase 2 Final Technical Report -- Antarctic penguins}

\documentclass{article}
\pdfpagewidth=8.5in
\pdfpageheight=11in

\usepackage{ijcai26}
\usepackage{times}
\usepackage{soul}
\usepackage{url}
\usepackage[hidelinks]{hyperref}
\usepackage[utf8]{inputenc}
\usepackage[small]{caption}
\usepackage{graphicx}
\usepackage{amsmath}
\usepackage{booktabs}
\usepackage{array}

\urlstyle{same}

% PDF Info Is REQUIRED. (\pdfinfo is a pdfTeX primitive; emit the same
% metadata via a \special under XeTeX so the PDF carries it either way.)
\ifx\pdfinfo\undefined
  \AtBeginDocument{\special{pdf:docinfo << /TemplateVersion (IJCAI.2026.0) >>}}
\else
  \pdfinfo{
  /TemplateVersion (IJCAI.2026.0)
  }
\fi

\title{UrbanTripV6: A Fault-Tolerant Neuro-Symbolic System for\\
ChinaTravel Phase 2}

\author{
Gangyi Zhang$^1$\thanks{Team leader}\and
Liangfei Jin$^{2}$\and
Junyi Liu$^{2}$\and
Qinzhe Wang$^{2}$\And
Zhuolin Li$^{2}$\\
\affiliations
$^{1}$National University of Singapore\\
$^{2}$University of New South Wales\\
\emails
zhanggangyi1224@gmail.com
}

\begin{document}

 
\maketitle
 
\begin{abstract}
This paper describes the phase-2 entry, UrbanTripV6, from team Antarctic penguins for the ChinaTravel Travel Planning Challenge~\cite{tpc,chinatravel}. The submitted configuration is designated as tpcagent. The Qwen3.6-27B checkpoint from the Qwen3 family~\cite{qwen} is used only to translate natural language into DSL. Ranging from extraction and canonicalization through grounding, planning, and repairing, all downstream steps did not use any additional embedding-based parser.

UrbanTripV6 operates by first ranking candidate bundles, then scheduling activities via memoized DFS. Its heuristic design incorporates geographic anchors, budget bounds, adaptive candidate retrieval, and a restricted intracity transit index. With support for resumable outputs, up to eight independent workers share a single SGLang endpoint. The workers handle inherited budget accounting,  checking generated constraints, deterministic retries and a fallback path derived from database.

We ran a test of 100 queries for internal validation on an oracle-stripped familiarization simulation, executed on the formal serving stack (SGLang 0.5.10.post1, two A800 GPUs, tensor parallelism two). In this development-time measurement, the submitted configuration achieved an Overall score of 97.70, an FPR of 98, and a C-LPR of 99.82. Note that this score is not an organizer-held-out evaluation.


\end{abstract}
 
\section{Task and System Overview}
\label{sec:scope}
 
ChinaTravel asks a system to turn a natural-language request into a
database-grounded multi-day itinerary that satisfies schema constraints,
commonsense constraints, and personal constraints, while also optimizing
three soft objectives. The official scoring formula is
\begin{equation*}
\begin{split}
S={}&.10\,\mathrm{EPR_{mic}}+.10\,\mathrm{EPR_{mac}}
   +.25\,\mathrm{C\text{-}LPR}+.40\,\mathrm{FPR}\\
  &+.05\,(\mathrm{DAV}+\mathrm{ATT}+\mathrm{DDR}).
\end{split}
\end{equation*}
Because FPR and C-LPR carry most of the weight, feasibility is effectively
the top priority.
 
\paragraph{Architecture.}
UrbanTripV6 is built around the four-stage neuro-symbolic pipeline shown in
Table~\ref{tab:pipeline}. Qwen's role is confined to interpreting the request -- it neither produces
the itinerary nor drives ranking. The symbolic stages that follow stabilize
the model's output, build a database-grounded plan, and accept only
independently verified changes, all wrapped in a crash-isolated execution
layer managing shared deadlines, resumable outputs, and a database-derived
fallback.
 
\begin{table*}[t]
\centering
\footnotesize
\setlength{\tabcolsep}{4.5pt}
\begin{tabular}{@{}p{3.0cm}p{4.7cm}p{8.1cm}@{}}
\toprule
Stage & Transformation & Main mechanisms \\
\midrule
Request interpretation & Natural language $\rightarrow$ raw DSL & Two-step Qwen translation and robust extraction of the generated constraint list. \\
Constraint stabilization & Raw DSL $\rightarrow$ grounded state & Canonicalization, executable and AST checks, targeted reflection, span-based coverage recovery, and database grounding. \\
Symbolic planning & Constraints $\rightarrow$ itinerary & Bundle enumeration, geographic and budget bounds, memoized DFS, and transit-aware ordering. \\
Verified refinement & Itinerary $\rightarrow$ final plan & Four feasibility repairs and twelve soft-metric stages, gated by the generated DSL. \\
\bottomrule
\end{tabular}
\caption{Four-stage architecture of UrbanTripV6.}
\label{tab:pipeline}
\end{table*}
 
\section{Constraint Compilation}
 
\paragraph{Translation and Isolation.}
Phase 2 works with English requests via the \texttt{nl2sl\_hybrid\_en}
translation pipeline. Qwen3.6-27B first turns the request into a structured
list of requirements, then converts that list into executable, Python-like
DSL blocks, so that the second call concentrates purely on DSL syntax
rather than on how the itinerary should be built. The model client issues OpenAI-compatible
requests with zero temperature and a fixed base seed. Grammar-constrained
decoding is kept, but aligned with the prompt: list-shaped calls request
one exact object shape and constrain decoding to that JSON schema, since a
generic \texttt{json\_object} grammar forbids the very array the prompt
asks for. A ladder of retries falls back to \texttt{json\_object} and then
to a bare payload whenever the endpoint rejects an optional field. Only the generated program is
passed downstream -- neither the model's natural-language reasoning trace
nor any itinerary it might propose ever reaches the search stage. The
prompts demand one-to-one coverage of explicit requirement clauses and
disallow room or bed constraints unless the request actually licenses them;
a deterministic filter, gated on the NL text, strips out any such
ungrounded constraint that still slips through. All four
\texttt{hard\_logic*} reference fields are stripped before translation and
planning. Disabling the thinking mode and capping output at 8192 tokens
both help cut down truncation-related variance. Every previously unseen
request is translated on the fly.
 
\paragraph{Robust Extraction and Canonicalization.}
Because a strict \texttt{json\_object} grammar can wrap the requested
constraint array inside an outer object, the extractor repairs and parses
the full response, unwraps any list-valued field or object-of-strings, and
otherwise scans every balanced bracket span, keeping the longest candidate
that parses as a list of strings so that an inner list literal is never
mistaken for the constraint list. The extracted DSL is then normalized for aliases, quote
variants, set-variable naming, budget accumulators, transport guards, and
database type spellings. Runtime checks on example plans and an AST checker can trigger up to five
general reflection calls; a deterministic disjunction-gap detector may
issue up to two targeted ones. A final mechanical coverage pass then treats the request text as the
authority for the clause families whose phrasing is templated. Each
generated constraint is cross-checked against its source clause on
polarity, window endpoints, and budget scope, and the canonical form is
synthesized whenever they disagree or the clause has no counterpart at
all: requirements rendered as prohibitions are flipped back, dropped
exclusion lists and type-inclusion lists are reinstalled, and a
type-guarded cost accumulator is replaced by the unguarded one its budget
noun implies. This recovers explicit budgets, transport modes, room
counts, POI types, and time windows. Checker feedback stays localized: a reflection call gets the specific
syntax or value diagnosis that failed, and disjunction repair fires only
when a required OR branch is missing. For clear-cut entity--category
mismatches, required POI names are recategorized only when every name is
absent from the generated category and present in exactly one other. After
this deterministic cleanup the benchmark's own UID, trip length, party
size, and city list are written back in. What comes out the other end is a stable constraint program,
rather than a loose, free-form semantic representation.
 
\paragraph{Literal DSL Parsing.}
The planner pulls a search state out of the DSL using block-local regular
expressions, scanners for quoted and set literals, named-set inlining, and
\texttt{ast.literal\_eval}. It handles both the set-algebra dialect and the
membership/flag dialect that Qwen tends to produce. For a disjunctive
request, each branch is searched on its own; candidates are then checked
against the original, unbranched generated DSL and the first to pass is
kept, falling back to whichever branch succeeded at the branch level. The
quoted-literal scanner keeps bracket characters intact inside strings and
handles escaped quotes before \texttt{ast.literal\_eval}; named set
variables are inlined before polarity, count, transport, hotel, meal,
attraction, and budget states get installed. No embedding similarity is
involved anywhere: everything after generation is grounded in literal text
spans, and reference constraints are never loaded at inference time.
 
\section{Constraint-Grounded Symbolic Planning}
 
\paragraph{Bundle Enumeration.}
UrbanTripV6 enumerates bounded triples consisting of outbound transport,
return transport, and accommodation. Caps on the number of intercity and
accommodation components, hotel property requirements, and single-day
chronology all act as hard filters on these bundles. A direction-specific
mode filter that would otherwise eliminate every available leg in one
direction is independently relaxed back to the raw timetable, so that
best-effort planning remains possible. Required POI names feed into hotel
scoring, while required-POI constraints inform the reachability checks.
Intercity and lodging costs, plus a floor for required-restaurant spend,
form a soft overall-budget preference: bundles exceeding it are dropped
only if another bundle satisfies it. The first pass also applies
required-POI reachability and a metro-only late-arrival heuristic; if that
empties the candidate set, a bounded second pass bypasses those two
POI-side heuristics while all generated-DSL validation stays in force.
Surviving bundles are ranked on normalized intercity cost, hotel cost, an
estimated geographic route, a hard-feature bonus, and travel-day utility,
the last rewarding early arrivals and late departures so the first and last
days keep usable activity time. Cheap and timing-extreme options are merged
before truncation, and the submitted configuration uses no model-based
reranking of bundles.
 
\paragraph{Memoized DFS.}
For each ranked bundle, a depth-first planner lays out meals, attractions,
hotel stays, and the grounded transport chains connecting them. A
coarse-grained memo signature made up of
$(\mathrm{day},\mathrm{time},\mathrm{position},\mathrm{uncovered\ requirements})$
is used to suppress repeated states. Whenever required items are still
outstanding, the normal 50-candidate search width is narrowed down to just
eight high-priority candidates. Each expansion pulls route mode, duration, distance, cost, and endpoints
from the travel environment; candidates are filtered before recursion by
opening hours, remaining budget, still-needed POI types, and allowed modes.
A hard-constraint-aware fallback repair and a budget-drop pass preserve a
usable best-effort plan when the search cannot finish in time. Given a
fixed compiled DSL, this ordering is entirely rule-driven.
 
\paragraph{Slack-Aware Deterministic Ranking.}
Budget slack and time slack feed into the serial DFS through
pressure-dependent ordering of candidates. Budget pressure is defined as
spend divided by a binding cap; time pressure combines trip progress with
time remaining until the return leg; pending pressure is the fraction of
required items still uncovered. Each candidate POI gets a normalized score
of the form
\begin{equation*}
\begin{split}
h={}&w_c\hat c+w_d\hat d+w_x\hat x+w_s\hat s\\
   &+w_p\hat p+w_m\hat m-w_g g+w_o b,
\end{split}
\end{equation*}
where $c,d,x,s,p,m$ stand for route cost, distance, transit duration, base
segment score, POI price, and closing-time margin, respectively; $g$
captures required-item coverage gain and $b$ is an ordering penalty. Lower
$h$ is better. The deterministic weighting leans on price/route pressure as budgets
tighten and on coverage/time-margin pressure as the trip nears its end,
with the minimum transit-time weight fixed at 3.0. Direct budget and
time-gap checks also shape later meal and attraction insertions.
 
\paragraph{Restricted Precomputed Segment Signal.}
An offline step precomputes mode-specific walk, metro, and taxi route
records for covered POI pairs, packaged as JSONL and loaded at inference
time keyed by city, ordered endpoints, and mode. Segment-based search is
disabled for hotels and intercity travel
(\texttt{use\_segments=False}); in its place a lazy
\texttt{SegmentIndex(build\_tfidf=False)} returns, for intracity lookups,
the minimum duration across the available mode variants, falling back to a
geodesic taxi-time estimate with a coverage penalty for uncovered pairs.
Segment data therefore influences only POI ranking and required-POI
reachability, never the emitted transport chains, which stay grounded by
the travel environment. Because no TF--IDF vectors are built at inference
time, the table functions purely as routing metadata rather than as
semantic retrieval.
 
\paragraph{Planner-Side Completion.}
V6 tracks the strongest candidates seen during search on logical,
commonsense, and schema criteria. Before returning, it takes the best
nonempty candidate or else assembles a database skeleton, then applies
output normalization, missing-meal insertion, required-evening-POI repair,
and DAV-oriented attraction insertion -- all ahead of the separate
refinement battery described next.
 
\section{Verified Repair and Robust Execution}
 
\paragraph{Phase-2 Correctness Adaptations.}
Results that lack a matching generated \texttt{hard\_logic\_py} field skip
refinement entirely. Plans emitted from the best-effort tail after a
planner timeout are admitted deliberately: that tail is exactly where the
search has already relaxed its own pruning, so those plans need repair
most, and the battery's own deadline guard bounds the extra time.
The first four repair stages target the tightened Phase-2 feasibility
checks. \texttt{dedupmeal} removes duplicate meal types, keeping entries
that are in-window, priced, and earliest in that order, then rethreads the
timeline. \texttt{deoverlap} shifts an overlapping activity to the previous
end time, preserving its duration when that still fits before midnight, and
rewrites the inbound chain without moving its endpoints; chronology is
repaired within a day while spatial continuity is tracked across days.
\texttt{fixspace} detects incorrect route endpoints and rebuilds the chain
through the travel environment with database-consistent mode, cost,
distance, duration, and timestamps. Finally, \texttt{mustpoi} parses the several set and literal
dialects that show up in the generated DSL and applies grouped repairs
across hotels, meals, attractions, intercity travel, and budget; its main
insertion logic respects whatever mode restrictions were parsed out.
Repair is deliberately local rather than a rebuild: a required meal window
is reached by absorbing slack in a later transport chain, a banned
intracity mode is re-routed leg by leg, and a paid meal that no cheaper
venue can replace is deleted with the following chain restitched.
Multi-item targets are repaired and gated as a group, since fixing just one
item in isolation may not actually improve satisfaction of a set-valued
constraint. Each stage applies its own acceptance test: \texttt{dedupmeal} needs a
non-regressing commonsense micro-score (strictly improving while the
candidate still fails commonsense); \texttt{fixspace} needs the commonsense
verdict to flip fail-to-pass; \texttt{deoverlap} needs the repaired plan to
pass outright; \texttt{mustpoi} needs a strict increase in aggregate
generated-constraint satisfaction. None may lose an existing pass.
 
\paragraph{Conditional Soft Refinement.}
Twelve further stages add, move, or reorder meals and attractions, ordered
so that the meal-completeness stages run before the costlier attraction
ones: \texttt{endday}, \texttt{gapmeal}, \texttt{travelday},
\texttt{bfstack}, \texttt{att}, \texttt{gapattr}, \texttt{endattr},
\texttt{dav2}, \texttt{fillerswap}, \texttt{seqswap}, \texttt{lateattr},
and \texttt{attdilute}. Each edit is kept only if it clears a stage-specific
non-regression check and actually improves the target metric.
\texttt{bfstack} allows a single hotel breakfast but blocks stacking
multiple breakfasts. Every gate in this battery relies on the generated
DSL rather than reference constraints. Functionally, these soft stages
cover meal completeness, attraction density, use of travel days, and local
sequence or filler swaps. One attraction-insertion stage restricts itself
to free POIs and walking whenever overall or local-transport budget slack
is running low. The whole battery stops launching new stages once fewer
than five seconds remain. Each stage operates on a copy of the plan and
simply returns its input unchanged if it's inapplicable or hits an error,
which limits how far any single regression can cascade.
 
\paragraph{Parallel and Failure-Safe Execution.}
The supervisor picks a worker count of
$N=\max(1,\min(8,\lfloor\mathrm{CPU}/2\rfloor))$ and shards queries by UID
index modulo $N$. Workers are subprocesses of the same \texttt{tpcagent}
invocation rather than parallel branches of the symbolic search: each keeps
its own planner state, works its shard sequentially with DFS serial within
a query, and reuses a single model client, while per-worker BLAS thread
caps keep the CPU from being oversubscribed. A daemon watchdog bounds the
wait and triggers the fallback at the emission deadline. Any existing
per-UID JSON file is a resume marker; a restarted worker inherits the
original start time and the shared soft budget of 17,400 seconds (4h50m).
After two worker crashes on one UID the system is forced onto the
deterministic database-fallback path, and if a worker exits abnormally the
supervisor may run a serial pass over unfinished queries before the hard
stop. All workers share one SGLang server, so sharding buys fault isolation
and concurrent submission rather than linear throughput.
 
\paragraph{Generated-Constraint Self-Check and Retry.}
After each attempt the runner checks schema compliance, commonsense, and
the generated DSL, treating a list of fewer than seven constraints as too
short. Multi-worker execution allows at most three attempts per query: a
failing or too-short attempt is retried only if at least 2,400 seconds
remain on the shared clock, clearing the translation cache, bumping the
base seed by 1,000, and capping each retry at 600 seconds. Candidates are
ranked lexicographically by full-pass status, length threshold, commonsense
pass, and satisfied-constraint count. Reference constraints and the
official Overall score play no role in this selection.
 
\paragraph{Budgeting and Deadlines.}
Single-worker runs use a 170-second outer timeout, raised to 1,200 seconds
under the supervisor. The wrapper works on a 290-second budget, sets V6's
external timeout to 245 seconds and reserves 60 for refinement -- roughly a
185-second core cutoff from the start of translation -- while the emission
watchdog fires near 275 seconds; later-triggering supervisor guards handle
defensive restarts and hard termination.
 
\section{Phase-2 Reproduction and Results}
 
\paragraph{Runtime Configuration.}
The target environment is Debian 12 with Python 3.12; SGLang 0.5.10.post1
serves Qwen3.6-27B over the OpenAI-compatible
\texttt{/v1/chat/completions} endpoint on local port 30000, on two A800
GPUs with tensor parallelism two and thinking disabled. The submitted
configuration is \texttt{tpcagent} with English input, served model name
\texttt{Qwen3.6-27B}, and local API key \texttt{EMPTY}. Evaluation runs
\texttt{solve\_script\_with\_harness.py} for a chosen method and split,
writing each itinerary to \path{results/<method>/<uid>.json}; oracle-based
aggregate scoring then consumes that directory.
 
\paragraph{Oracle-Free Simulation and Deliverable.}
The public \texttt{phase2\_heldout\_sim} split holds the same 100 requests
as \texttt{phase2\_familiar} with the reference-constraint fields stripped.
We run inference on this oracle-free split with a 100-query limit, save one
JSON file per UID, and then score the completed directory against the
familiarization references without rerunning the planner; the formal
held-out run follows the same output contract on organizer-provided
requests. FPR is the intersection of schema, commonsense, and logical pass
IDs; failing plans are assigned zero DAV/ATT/DDR before the score of
Section~\ref{sec:scope} is computed.
 
\paragraph{Packaged-Configuration Rehearsal.}
For the packaged, retry-enabled configuration (archive MD5 prefix
\texttt{22bcd0d6}), our internally recorded 100-query result was FPR 98,
C-LPR 99.82, both EPR terms at 100, and Overall 97.70. Inference used the
oracle-stripped familiarization simulation on the serving stack described
above, with translations regenerated from scratch; reference constraints
were only reintroduced for aggregate scoring after every plan had already
been written out. Successive stages of the constraint verifier and the
repair battery moved this measurement along 70.5, 83.2, 93.8, and 97.7
Overall (FPR 52, 73, 91, 98). These are development-time measurements, not
organizer-held-out scores -- the latter simply wasn't available at
submission time.
 
\section{Limitations and Conclusion}
 
Translation remains the biggest risk to generalization: clause-anchored
repair recovers the templated families, but cannot guarantee that Qwen
preserved every semantic clause of an unseen free-form request. The
segment graph is deliberately incomplete, so uncovered routes fall back on
a geographic estimate. Compound failures needing a whole day rebuilt may
exceed the generic battery, the residual familiarization failures come
from run-to-run translation variance on borderline queries and from a few
geometrically over-constrained instances, and all workers compete for one
served model. Overall, the design keeps language interpretation, literal
parsing and canonicalization, pressure-aware symbolic search,
constraint-gated repair, and crash-isolated execution as distinct
components; together they make up the Phase-2 system submitted here.
 
\begin{thebibliography}{}
\bibitem[\protect\citeauthoryear{LAMDA-NeSy}{2026}]{tpc}
LAMDA-NeSy.
\newblock Travel Planning Challenge (TPC) @ IJCAI-ECAI 2026.
\newblock \url{https://chinatravel-competition.github.io/IJCAI2026/}, 2026.
 
\bibitem[\protect\citeauthoryear{Qwen Team}{2025}]{qwen}
Qwen Team.
\newblock Qwen3 Technical Report.
\newblock arXiv:2505.09388, 2025.
 
\bibitem[\protect\citeauthoryear{Shao et al.}{2026}]{chinatravel}
Jie-Jing Shao, Bo-Wen Zhang, Xiao-Wen Yang, \emph{et al.}
\newblock ChinaTravel: An Open-Ended Travel Planning Benchmark with
Compositional Constraint Validation for Language Agents.
\newblock In \emph{ICLR}, arXiv:2412.13682, 2026.
\end{thebibliography}

\end{document}
